% !TEX root = ../main.tex

\chapter{绪论}

\section{研究背景与意义}

人形机器人兼具类人的运动形态与任务交互能力，被普遍认为是具身智能系统迈向通用场景的重要平台。与轮式机器人相比，人形机器人能够适应楼梯、狭窄通道和以人类尺度设计的工作环境；与固定机械臂相比，人形机器人又天然具备移动能力和更大的操作工作空间，因此在家庭服务、工业巡检、灾害救援以及危险环境替代作业等场景中展现出显著潜力。

然而，人形机器人要真正进入开放环境并执行复杂任务，首先必须解决全身运动控制问题。该问题同时面临高自由度、动力学强耦合、接触切换频繁以及系统欠驱动等多重挑战。下肢需要维持稳定步态和抗扰平衡，上肢需要完成灵巧操作和精确跟踪，而二者在质心变化、关节力矩分配和接触约束方面又彼此影响。传统依赖模型推导和人工设计的控制方法虽然具有较好的可解释性，但面对复杂场景、多任务切换以及高自由度形态时，往往难以兼顾稳定性、灵活性与开发效率。

在这一背景下，“人在回路”的研究范式为人形机器人控制提供了一条兼顾数据效率与任务落地的新路径。通过遥操作系统获取高质量人体动作数据，可以显著降低复杂技能构建的门槛；通过动作重定向与控制接口设计，可以将人体示教经验迁移到机器人平台；进一步结合强化学习和感知闭环控制，则能够推动系统从“人控执行”发展到“自主执行”。因此，围绕遥操作数据采集、全身控制策略和任务级闭环验证建立统一技术链路，不仅具有重要理论意义，也对推动人形机器人走向真实应用具有工程价值。

\section{国内外研究现状}

\subsection{遥操作与数据采集}

遥操作是连接人类动作能力与机器人执行能力的重要桥梁，也是当前构建高质量具身数据集的重要手段。已有工作通常通过 VR 设备、惯性传感器、光学动捕系统或力反馈装置采集人体运动信息，并将其映射到机械臂、双臂平台或移动操作机器人。此类方法在示教采集、快速验证和远程控制方面效果显著，但在面向人形机器人的全身场景时仍面临若干问题。

首先，不同采集设备的精度、延迟和部署成本差异较大，单一数据源往往难以同时覆盖高精度研究、日常采集与复杂交互任务三类需求。其次，人体与机器人之间普遍存在形态差异，包括肢体长度不一致、关节自由度不同以及执行约束不一致等，若缺乏统一的重定向机制，直接映射很容易造成动作失真、轨迹突变或关节超限。再次，当前不少系统更侧重上肢操作或局部动作采集，对移动与操作耦合的全身行为支持仍不充分。

因此，面向人形机器人的遥操作系统需要同时考虑多源数据接入、形态差异建模与机器人可执行性约束。本文针对上述问题构建多源人体运动捕捉体系，并以 GMR 为核心建立统一动作重定向框架，提升从人体动作到机器人动作的可迁移性与一致性。

\subsection{通用全身运动控制}

全身运动控制的目标是在满足接触约束和稳定性要求的前提下，使机器人协调完成移动、平衡和操作等行为。传统方法多依赖动力学模型、优化求解器以及分层控制架构，通过逆动力学、QP 优化或模型预测控制等方法实现任务优先级分配。这类方法在规则场景中具有较强稳定性，但在复杂接触、多样动作和高维策略搜索方面通常需要大量模型调参与任务定制。

随着强化学习在机器人领域的广泛应用，越来越多工作尝试利用仿真训练获得鲁棒步态和复杂运动技能。强化学习方法能够在大规模并行环境中自动搜索控制策略，尤其适合下肢步态生成和抗扰动控制。然而，完全端到端的全身策略训练通常需要更高训练成本，也对奖励设计、状态表征和仿真真实性提出更高要求。此外，端到端全身策略在上肢精细操作方面不一定具备足够优势，且策略解释性和工程调试难度较高。

因此，一种兼顾稳定移动与灵巧操作的可行路径是采用上下肢解耦的控制设计思路：下肢以强化学习获得通用 Locomotion 能力，上肢以轨迹跟踪和逆运动学保证操作精度，再通过统一接口实现全身协调。本文即基于这一思路设计通用全身运动控制框架，在系统层面平衡灵活性、鲁棒性和可部署性。

\subsection{感知--控制闭环与自主任务}

仅具备离线控制能力的人形机器人仍难以在开放环境中稳定执行任务。为了实现真正的自主作业，系统必须能够感知外界状态、理解任务目标、实时生成控制指令并根据反馈持续修正行为。近年来，视觉感知、语音理解、大模型规划与机器人控制的结合推动了感知--控制闭环系统的发展，使机器人逐步具备从目标识别到任务执行的端到端能力。

在动态交互任务中，感知模块通常需要以较高频率完成目标检测、状态估计与轨迹预测，控制模块则需要在严格时延约束下生成全身动作；在服务型任务中，高层规划器需要将自然语言指令分解为结构化子任务，视觉模块需要输出可用于抓取和导航的目标位姿，而底层控制器则负责在动态环境中完成移动、抓取和递送。尽管相关研究不断推进，但感知输出与全身控制接口之间仍常存在抽象层级不一致、执行链路不统一和系统复用性不足的问题。

本文关注的不是单一感知算法本身，而是感知模块与全身控制系统之间的接口构建与任务闭环验证。通过在 CyboRacket 和 OpenClaw 两类任务中进行系统级集成，本文验证了从感知结果到全身动作执行的统一控制链路。

\section{本文主要研究内容}

本文围绕“数据采集--控制策略--任务执行”三层主线展开研究，目标是建立一套可支撑人形机器人全身运动控制与自主任务执行的统一框架。具体而言，本文的主要工作包括以下三个方面。

第一，面向遥操作和具身数据采集需求，构建多源异构人体运动捕捉体系，涵盖 Noitom 惯性动捕、Nokov 光学动捕和基于 SONIC 的仿真交互数据。针对人体与机器人之间的形态差异，设计基于 GMR 的动作重定向方法，将人体动作转换为机器人可执行的关节轨迹或末端目标。

第二，针对人形机器人全身控制中移动与操作相互耦合的问题，提出上下肢解耦的通用控制架构。下肢部分基于强化学习训练 Locomotion 策略，以获得稳定的基础移动能力；上肢部分构建面向遥操作和自动规划结果的统一轨迹接口，并结合逆运动学实现高保真跟踪；同时引入全身协调机制，用于处理上肢大范围动作对下肢稳定性的影响。

第三，围绕感知驱动的任务级应用，构建感知--控制闭环系统，并在动态球拍任务与语音指令服务任务中开展验证。通过将视觉感知、任务规划与全身控制接口统一起来，本文展示了所提方法在动态目标交互和服务型具身任务中的应用潜力。

\section{论文组织结构}

全文共分为五章，并附录相关硬件与训练配置说明。

第 1 章为绪论，介绍课题研究背景、相关研究现状、本文主要研究内容以及论文整体结构。

第 2 章讨论多源遥操作数据采集与重定向问题，重点介绍多源运动捕捉体系、动作重定向建模方法以及重定向精度评估思路。

第 3 章介绍解耦式通用全身运动控制策略，重点阐述上下肢解耦架构、下肢强化学习 Locomotion 策略、上肢轨迹跟踪接口以及 Sim-to-Real 真机验证方法。

第 4 章给出感知--控制闭环自主任务执行的系统实现，选取 CyboRacket 和 OpenClaw 两个任务作为验证案例，展示感知、规划与全身控制的系统级融合方式。

第 5 章总结全文工作，归纳主要创新点，并对未来研究方向进行展望。
